\documentclass{article}
\usepackage{float}
\usepackage{svg}
\usepackage{amsmath}
\usepackage[utf8]{inputenc}
\usepackage[french]{babel}
\usepackage{graphicx}
\usepackage[T1]{fontenc}
\usepackage{hyperref}
\usepackage{changepage}
\hypersetup{
	colorlinks,
	citecolor=black,
	filecolor=black,
	linkcolor=black,
	urlcolor=black
}
%\graphicspath{ressources/}

\usepackage{blindtext}

\usepackage{subfiles} % Best loaded last in the preamble

\title{Titre}
\author{ }
\date{ }

\begin{document}
\begin{titlepage}
	\begin{center}

		{\Large Université de Mons}\\[1ex]
		{\Large Faculté Des Sciences}\\[1ex]
		{\Large Département d'Informatique}\\[1ex]

		\newcommand{\HRule}{\rule{\linewidth}{0.3mm}}
		% Title
		\HRule \\[0.3cm]
		{ \LARGE \bfseries Projet réseau - Selective-repeat \\[0.3cm]}
		%{ \LARGE \bfseries Rapport d'implémentation  \\[0.1cm]} % Commenter si pas besoin
		\HRule \\[1.5cm]

		% Author and supervisor
		\begin{minipage}[t]{0.45\textwidth}
			\begin{flushleft} \large
				\emph{Professeur:}\\
				Bruno \textsc{Quoitin}\\
				\emph{Assistants:}\\
				Jeremy \textsc{Dubrulle}\\
			\end{flushleft}
		\end{minipage}
		\begin{minipage}[t]{0.45\textwidth}
			\begin{flushright} \large
				\emph{Auteurs:} \\
				Théo \textsc{Godin} \\
				François \textsc{Vion} \\
			\end{flushright}
		\end{minipage}\\[2ex]

		\vfill

		% Bottom of the page
		\begin{center}
			\begin{tabular}[t]{c c c}
				\includegraphics[height=1.5cm]{ressources/logoumons.jpg} &
				\includegraphics[height=1.5cm]{ressources/logofs.jpg} &
			\end{tabular}
		\end{center}~\\

		{\large Année académique 2021-2022}

	\end{center}
\end{titlepage}

\tableofcontents
\newpage

\pagenumbering{arabic}

\section{Applications}

Nous avons décidé de suivre l'exemple ping-pong et de créer deux classes application, une qui envoie les données et une autre qui les reçcoit.
L'app sender recoit les données à envoyer en paramètre de son constructeur et passe ensuite ces données à la méthode sendMessage du protocol. 
Tandis que l'app receiver, au lancement de l'application, ne fait qu'écouter.

\section{Selective-repeat}

\subsection{Envoi d'un message}

Un message peut être envoyé par une application via la méthode sendMessage de notre prtocol. Cette méthode va séparer ce message en plus petits messages de 8 caractères qui vont pouvoir ensuite être envoyés en pipelining.

\subsection{Pipelining}

La méthode pipeliningSend s'occupe d'envoyer des segments tant que la window n'a pas atteint sa taille maximale.

\subsection{Envoi d'un segment}

Afin d'envoyer un segment, la méthode send envoie ce segment à la destination, l'ajoute à l'arrayList représentant la window et finalement crée et démarre un timer afin de gérer les timeout pour ce segment.

\subsection{Timeout}

Si un timer arrive à sa fin, le segment associé est ré-envoyé et un nouveau timer est créé.

\end{document}
